\documentclass[11pt,a4paper]{article}

%====================================================
% PACKAGES
%====================================================
\usepackage[utf8]{inputenc}
\usepackage[T1]{fontenc}
\usepackage[french]{babel}

\usepackage[a4paper,margin=2cm]{geometry}

\usepackage{amsmath,amssymb}
\usepackage{lmodern}

\usepackage{fancyhdr}
\usepackage{enumitem}
\usepackage{xcolor}

%====================================================
% MISE EN PAGE
%====================================================
\pagestyle{fancy}
\fancyhf{}

\fancyhead[L]{\textbf{Algèbre linéaire}}
\fancyhead[C]{\textbf{TD -- Chapitre 8}}
\fancyhead[R]{\thepage}

\setlength{\headheight}{14pt}

%====================================================
% TITRE
%====================================================
\begin{document}

\begin{center}
{\Large \textbf{Classe préparatoire -- Mathématiques}}

\vspace{0.3cm}

{\Huge \textbf{TD -- Applications linéaires}}

\vspace{0.3cm}

{\large Chapitre 8}

\vspace{0.5cm}

\textbf{Nom :} \underline{\hspace{5cm}} \hfill
\textbf{Date :} \underline{\hspace{3cm}}
\end{center}

\vspace{0.5cm}

\hrule
\vspace{0.3cm}

\textbf{Objectifs :}
\begin{itemize}
\item Maîtriser la notion d'application linéaire
\item Étudier noyau, image et rang
\item Comprendre les structures fondamentales
\item Se préparer aux oraux de concours
\end{itemize}

\vspace{0.3cm}
\hrule

\vspace{0.8cm}


\section{Exercices d'application directe}

\textbf{Exercice 1 [APP $\star$] — Vérification de linéarité}

On considère l'application $u : \mathbb{R}^2 \to \mathbb{R}^2$ définie par :
\[
u(x,y) = (2x - y,\, x + 3y)
\]

\begin{enumerate}
\item Vérifier que $u(0,0) = (0,0)$.
\item Vérifier que pour tous $(x_1,y_1)$ et $(x_2,y_2)$ :
\[
u((x_1,y_1)+(x_2,y_2)) = u(x_1,y_1) + u(x_2,y_2)
\]
\item Vérifier que pour tout $\lambda \in \mathbb{R}$ :
\[
u(\lambda(x,y)) = \lambda u(x,y)
\]
\item Conclure que $u$ est une application linéaire.
\end{enumerate}

---

\textbf{Exercice 2 [APP $\star$] — Contre-exemple}

On considère l'application $f : \mathbb{R}^2 \to \mathbb{R}$ définie par :
\[
f(x,y) = x^2 + y
\]

\begin{enumerate}
\item Calculer $f(x_1+x_2, y_1+y_2)$.
\item Comparer avec $f(x_1,y_1) + f(x_2,y_2)$.
\item Conclure que $f$ n'est pas linéaire.
\end{enumerate}

---

\textbf{Exercice 3 [APP $\star\star$] — Noyau}

On considère l'application $u : \mathbb{R}^2 \to \mathbb{R}$ définie par :
\[
u(x,y) = x - y
\]

\begin{enumerate}
\item Déterminer l'ensemble des vecteurs $(x,y)$ tels que $u(x,y)=0$.
\item Montrer que cet ensemble est une droite vectorielle.
\item Donner une base du noyau.
\item En déduire sa dimension.
\end{enumerate}

---

\textbf{Exercice 4 [APP $\star\star$] — Image}

On considère l'application $u : \mathbb{R}^3 \to \mathbb{R}^2$ :
\[
u(x,y,z) = (x+y,\, y+z)
\]

\begin{enumerate}
\item Calculer $u(1,0,0)$, $u(0,1,0)$ et $u(0,0,1)$.
\item Montrer que l'image est engendrée par ces vecteurs.
\item Déterminer une base de $Im(u)$.
\item En déduire le rang de $u$.
\end{enumerate}

---

\textbf{Exercice 5 [APP $\star\star$] — Injectivité et surjectivité}

On considère $u : \mathbb{R}^3 \to \mathbb{R}^2$ :
\[
u(x,y,z) = (x+y,\, y+z)
\]

\begin{enumerate}
\item Résoudre $u(x,y,z)=(0,0)$.
\item En déduire le noyau.
\item L'application est-elle injective ?
\item Comparer les dimensions de départ et d'arrivée.
\item L'application peut-elle être surjective ?
\item Conclure.
\end{enumerate}

---

\section{Exercices de méthode}

\textbf{Exercice 6 [METH $\star\star$] — Construction d’une application}

On considère la base canonique $(e_1,e_2)$ de $\mathbb{R}^2$.

On définit une application $u$ par :
\[
u(e_1) = (1,2), \quad u(e_2) = (0,1)
\]

\begin{enumerate}
\item Écrire un vecteur $(x,y)$ comme combinaison de $e_1$ et $e_2$.
\item En déduire l'expression de $u(x,y)$.
\item Montrer que l'application est linéaire.
\end{enumerate}

---

\textbf{Exercice 7 [METH $\star\star\star$] — Matrice d’une application}

On considère l'application $u : \mathbb{R}^2 \to \mathbb{R}^3$ :
\[
u(x,y) = (x+y,\, 2x-y,\, y)
\]

\begin{enumerate}
\item Calculer $u(e_1)$ et $u(e_2)$.
\item Écrire la matrice associée.
\item Vérifier que :
\[
u(x,y) = A \begin{pmatrix} x \\ y \end{pmatrix}
\]
\end{enumerate}

---

\textbf{Exercice 8 [METH $\star\star\star$] — Noyau via matrice}

On considère :
\[
A =
\begin{pmatrix}
1 & 2 \\
2 & 4
\end{pmatrix}
\]

\begin{enumerate}
\item Résoudre $AX = 0$.
\item Décrire le noyau.
\item Donner une base.
\item En déduire la dimension.
\end{enumerate}

---

\textbf{Exercice 9 [METH $\star\star\star$] — Image via matrice}

Même matrice que précédemment.

\begin{enumerate}
\item Identifier les colonnes de la matrice.
\item Montrer qu’elles sont liées.
\item En déduire une base de l’image.
\item Calculer le rang.
\end{enumerate}

---

\section{Exercices de démonstration}

\textbf{Exercice 10 [DEM $\star\star$] — Image de 0}

Soit $u$ une application linéaire.

\begin{enumerate}
\item Écrire $u(0)$ comme $u(0+0)$.
\item Utiliser la linéarité.
\item Conclure que $u(0)=0$.
\end{enumerate}

---

\textbf{Exercice 11 [DEM $\star\star\star$] — Noyau SEV}

Soit $u : E \to F$ linéaire.

\begin{enumerate}
\item Montrer que $0 \in Ker(u)$.
\item Montrer la stabilité par addition.
\item Montrer la stabilité par multiplication scalaire.
\item Conclure.
\end{enumerate}

---

\section{Exercices type ORAUX (Mines / Centrale)}

\textbf{Exercice 12 [ORAL $\star\star\star$]}

On considère :
\[
u(x,y,z) = (x+y, y+z, x+z)
\]

\begin{enumerate}
\item Déterminer le noyau.
\item Déterminer l’image.
\item Calculer le rang.
\item L’application est-elle injective ?
\item L’application est-elle surjective ?
\end{enumerate}

---

\textbf{Exercice 13 [ORAL $\star\star\star\star$] — Projecteur}

Soit $u$ telle que :
\[
u^2 = u
\]

\begin{enumerate}
\item Montrer que $Im(u)$ et $Ker(u)$ sont supplémentaires.
\item Montrer que :
\[
E = Ker(u) \oplus Im(u)
\]
\item Interpréter géométriquement.
\end{enumerate}

\section{Problèmes guidés (niveau concours)}

%====================================================
\textbf{Exercice 14 [PREP $\star\star\star$] — Étude complète d’une application}

On considère l’application $u : \mathbb{R}^3 \to \mathbb{R}^3$ définie par :
\[
u(x,y,z) = (x+y,\; y+z,\; x+z)
\]

\begin{enumerate}
\item Vérifier que $u$ est une application linéaire.
\item Déterminer le noyau de $u$.
\item En déduire la dimension du noyau.
\item Déterminer l’image de $u$.
\item En déduire le rang de $u$.
\item Vérifier le théorème du rang.
\item L’application est-elle injective ? surjective ?
\item Déterminer si $Im(u)$ est un plan, une droite ou tout $\mathbb{R}^3$.
\item Donner une interprétation géométrique de l'application.
\end{enumerate}

---

%====================================================
\textbf{Exercice 15 [PREP $\star\star\star$] — Construction d’une application}

On se place dans $\mathbb{R}^3$ muni de sa base canonique.

On cherche une application linéaire $u : \mathbb{R}^3 \to \mathbb{R}^2$ telle que :
\[
u(1,0,0) = (1,0), \quad
u(0,1,0) = (0,1), \quad
u(0,0,1) = (1,1)
\]

\begin{enumerate}
\item Montrer que ces conditions définissent une unique application linéaire en utilisant le fait qu'une application linéaire est entièrement déterminée par l'image d'une base.
\item Déterminer l’expression de $u(x,y,z)$.
\item Écrire la matrice de $u$.
\item Déterminer son noyau.
\item Déterminer son image.
\item Calculer le rang.
\end{enumerate}

---

%====================================================
\textbf{Exercice 16 [PREP $\star\star\star\star$] — Étude paramétrée}

On considère la famille d’applications :
\[
u_\lambda(x,y) = (x + \lambda y,\; y + \lambda x)
\]

\begin{enumerate}
\item Montrer que $u_\lambda$ est linéaire.
\item Écrire la matrice associée.
\item Calculer le déterminant.
\item Pour quelles valeurs de $\lambda$ l’application est-elle injective ?
\item Pour quelles valeurs de $\lambda$ est-elle surjective ?
\item Étudier le noyau pour $\lambda = 1$ et $\lambda = -1$.
\item En déduire le rang dans ces cas particuliers.
\end{enumerate}

---

%====================================================
\textbf{Exercice 17 [PREP $\star\star\star\star$] — Rang et dépendance}

On considère :
\[
u(x,y,z) = (x+y+z,\; x+2y+3z)
\]

\begin{enumerate}
\item Vérifier que $u$ est linéaire.
\item Calculer les images des vecteurs de la base canonique.
\item Montrer que ces vecteurs sont liés.
\item En déduire une base de l’image.
\item Calculer le rang.
\item Déterminer le noyau.
\item Vérifier le théorème du rang.
\end{enumerate}

---

%====================================================
\section{Problèmes de démonstration approfondie}

\textbf{Exercice 18 [DEM $\star\star\star$] — Image d’une famille génératrice}

Soit $u : E \to F$ une application linéaire.

Soit $(e_1,\dots,e_n)$ une famille génératrice de $E$.

\begin{enumerate}
\item Montrer que tout vecteur $x \in E$ s’écrit comme combinaison linéaire des $e_i$.
\item Appliquer $u$ à cette combinaison.
\item En déduire que $Im(u)$ est engendrée par $u(e_1),\dots,u(e_n)$.
\end{enumerate}

---

\textbf{Exercice 19 [DEM $\star\star\star\star$] — Image d’une famille libre}

Soit $u : E \to F$ une application linéaire injective.

\begin{enumerate}
\item Soit $(x_1,\dots,x_n)$ une famille libre.
\item Supposer une relation linéaire sur $(u(x_1),\dots,u(x_n))$.
\item Appliquer l’injectivité.
\item Conclure que la famille image est libre.
\end{enumerate}

---

\section{Exercices type ORAUX (Mines / Centrale)}

%====================================================
\textbf{Exercice 20 [ORAL $\star\star\star\star$] — Projecteur}

Soit $u : E \to E$ telle que :
\[
u^2 = u
\]

\begin{enumerate}
\item Montrer que pour tout $x$, $u(x)$ est un point fixe de $u$.
\item Montrer que $Im(u) \subset \{x \mid u(x)=x\}$.
\item Montrer que :
\[
E = Ker(u) + Im(u)
\]
\item Montrer que la somme est directe.
\item Conclure que $u$ est un projecteur.
\end{enumerate}

---

%====================================================
\textbf{Exercice 21 [ORAL $\star\star\star\star$] — Symétrie}

Soit $u : E \to E$ telle que :
\[
u^2 = Id
\]

\begin{enumerate}
\item Montrer que :
\[
(u-Id)(u+Id)=0
\]
\item Montrer que :
\[
E = Ker(u-Id) + Ker(u+Id)
\]
\item Montrer que la somme est directe.
\item Interpréter géométriquement.
\end{enumerate}

---

%====================================================
\textbf{Exercice 22 [ORAL $\star\star\star\star\star$] — Classification rang 1}

Soit $u : E \to F$ de rang 1.

\begin{enumerate}
\item Montrer que $Im(u)$ est une droite vectorielle.
\item Montrer qu’il existe $\varphi$ forme linéaire et $v \in F$ tel que :
\[
u(x) = \varphi(x) v
\]
\item Réciproquement, montrer que toute application de cette forme est de rang 1.
\end{enumerate}

---

%====================================================
\textbf{Exercice 23 [ORAL $\star\star\star\star\star$] — Théorème du rang (oral)}

Soit $u : E \to F$.

\begin{enumerate}
\item Construire une base adaptée du noyau.
\item Compléter en une base de $E$.
\item Étudier l’image des vecteurs restants.
\item En déduire :
\[
\dim(E) = \dim(Ker(u)) + rg(u)
\]
\end{enumerate}

---

%====================================================
\section{Problème final (type CCP / Mines)}
%====================================================

\textbf{Exercice 24 [PREP $\star\star\star\star\star$] — Étude globale}

On considère :
\[
u(x,y,z,t) = (x+y,\; y+z,\; z+t)
\]

\begin{enumerate}
\item Vérifier que $u$ est linéaire.
\item Déterminer le noyau.
\item Déterminer sa dimension.
\item Déterminer l’image.
\item En déduire le rang.
\item Vérifier le théorème du rang.
\item Étudier injectivité et surjectivité.
\item Donner une interprétation géométrique.
\end{enumerate}
\section{Problèmes avancés (niveau très élevé)}

%====================================================
\textbf{Exercice 25 [PREP $\star\star\star\star$] — Endomorphismes de rang 1}

Soit $u : E \to E$ une application linéaire de rang 1.

\begin{enumerate}
\item Montrer que $Im(u)$ est une droite vectorielle.
\item Soit $e \in Im(u)$ non nul. Montrer que pour tout $x \in E$, $u(x)$ est colinéaire à $e$.
\item En déduire qu’il existe une forme linéaire $\varphi$ telle que :
\[
u(x) = \varphi(x)\, e
\]
\item Montrer que toute application de cette forme est de rang 1 si $\varphi \neq 0$.
\item Étudier le noyau de $u$.
\item En déduire la dimension du noyau.
\end{enumerate}

---

%====================================================
\textbf{Exercice 26 [PREP $\star\star\star\star$] — Structure des projecteurs}

Soit $u : E \to E$ tel que $u^2 = u$.

\begin{enumerate}
\item Montrer que pour tout $x \in E$ :
\[
x = (x - u(x)) + u(x)
\]
\item Montrer que $x - u(x) \in Ker(u)$.
\item Montrer que $u(x) \in Im(u)$.
\item En déduire que :
\[
E = Ker(u) + Im(u)
\]
\item Montrer que cette somme est directe.
\item Conclure que $u$ est le projecteur sur $Im(u)$ parallèlement à $Ker(u)$.
\end{enumerate}

---

%====================================================
\textbf{Exercice 27 [PREP $\star\star\star\star$] — Décomposition canonique}

Soit $u : E \to E$.

On suppose que :
\[
E = Ker(u) \oplus F
\]

\begin{enumerate}
\item Montrer que la restriction de $u$ à $F$ est injective.
\item Montrer que $Im(u) = u(F)$.
\item En déduire que :
\[
\dim(Im(u)) = \dim(F)
\]
\item Retrouver le théorème du rang.
\end{enumerate}

---

%====================================================
\textbf{Exercice 28 [PREP $\star\star\star\star\star$] — Classification des endomorphismes}

Soit $u : E \to E$ tel que :
\[
u^2 = Id
\]

\begin{enumerate}
\item Montrer que :
\[
(u - Id)(u + Id) = 0
\]
\item Montrer que :
\[
Im(u - Id) \subset Ker(u + Id)
\]
\item Montrer que :
\[
E = Ker(u - Id) + Ker(u + Id)
\]
\item Montrer que la somme est directe.
\item Interpréter géométriquement cette décomposition.
\end{enumerate}

---

%====================================================
\section{Problèmes de type Mines / ENS}

%====================================================
\textbf{Exercice 29 [ORAL $\star\star\star\star\star$] — Théorème du rang (preuve avancée)}

Soit $u : E \to F$.

\begin{enumerate}
\item Soit $(e_1,\dots,e_k)$ une base du noyau.
\item Compléter cette famille en une base $(e_1,\dots,e_n)$ de $E$.
\item Montrer que les vecteurs $u(e_{k+1}),\dots,u(e_n)$ sont libres.
\item Montrer qu’ils engendrent l’image.
\item En déduire :
\[
\dim(E) = \dim(Ker(u)) + \dim(Im(u))
\]
\end{enumerate}

---

%====================================================
\textbf{Exercice 30 [ORAL $\star\star\star\star\star$] — Application nilpotente}

Soit $u : E \to E$ telle que :
\[
u^2 = 0
\]

\begin{enumerate}
\item Montrer que $Im(u) \subset Ker(u)$.
\item En déduire que :
\[
rg(u) \leq \dim(Ker(u))
\]
\item Montrer que :
\[
2\,rg(u) \leq \dim(E)
\]
\item Interpréter ce résultat.
\end{enumerate}

---

%====================================================
\textbf{Exercice 31 [ORAL $\star\star\star\star\star$] — Isomorphisme}

Soit $u : E \to F$.

\begin{enumerate}
\item Montrer que si $u$ est injective et $\dim(E)=\dim(F)$ alors $u$ est bijective.
\item Montrer que si $u$ est surjective et $\dim(E)=\dim(F)$ alors $u$ est bijective.
\item En déduire une caractérisation des isomorphismes.
\end{enumerate}

---

%====================================================
\section{Problème final (niveau élevé)}

\textbf{Exercice 32 [PREP $\star\star\star\star\star$] — Structure globale}

Soit $u : E \to E$ une application linéaire.

\begin{enumerate}
\item Montrer que :
\[
E = Ker(u) \oplus G
\]
pour un certain sous-espace $G$.

\item Montrer que la restriction de $u$ à $G$ est injective.

\item Montrer que :
\[
Im(u) = u(G)
\]

\item En déduire que $u$ induit un isomorphisme entre $G$ et $Im(u)$.

\item Interpréter géométriquement ce résultat.

\item En déduire une nouvelle preuve du théorème du rang.
\end{enumerate}
\section{Corrigé — Exercices d'application}

%====================================================
\textbf{Exercice 1 — Linéarité}

On considère :
\[
u(x,y) = (2x - y,\; x + 3y)
\]

\underline{1. Vérification de $u(0,0)$}

\[
u(0,0) = (0,0)
\]

---

\underline{2. Additivité}

Soient $(x_1,y_1)$ et $(x_2,y_2)$ :

\[
u((x_1,y_1)+(x_2,y_2))
= u(x_1+x_2,\; y_1+y_2)
\]

\[
= (2(x_1+x_2)-(y_1+y_2),\; (x_1+x_2)+3(y_1+y_2))
\]

\[
= (2x_1-y_1,\; x_1+3y_1) + (2x_2-y_2,\; x_2+3y_2)
\]

\[
= u(x_1,y_1) + u(x_2,y_2)
\]

---

\underline{3. Homogénéité}

\[
u(\lambda(x,y)) = u(\lambda x, \lambda y)
\]

\[
= (2\lambda x - \lambda y,\; \lambda x + 3\lambda y)
\]

\[
= \lambda (2x-y,\; x+3y)
\]

\[
= \lambda u(x,y)
\]

---

\underline{Conclusion}

$u$ est linéaire.

---

%====================================================
\textbf{Exercice 3 — Noyau}

\[
u(x,y) = x-y
\]

On cherche :
\[
x-y=0
\]

\[
x=y
\]

Donc :
\[
Ker(u) = \{(t,t),\; t \in \mathbb{R}\}
\]

---

\underline{Base}

\[
Ker(u) = Vect((1,1))
\]

---

\underline{Dimension}

\[
\dim(Ker(u)) = 1
\]

---

%====================================================
\textbf{Exercice 4 — Image}

\[
u(x,y,z) = (x+y,\; y+z)
\]

On calcule :

\[
u(1,0,0) = (1,0)
\]
\[
u(0,1,0) = (1,1)
\]
\[
u(0,0,1) = (0,1)
\]

---

On remarque :

\[
(1,1) = (1,0) + (0,1)
\]

Donc :

\[
Im(u) = Vect((1,0),(0,1)) = \mathbb{R}^2
\]

---

\underline{Conclusion}

\[
rg(u) = 2
\]

---

%====================================================
\textbf{Exercice 5 — Injectivité / Surjectivité}

Noyau :

\[
x+y=0,\quad y+z=0
\]

Donc :

\[
y = -x,\quad z = -y = x
\]

\[
Ker(u) = \{(x,-x,x)\}
\]

---

\underline{Dimension}

\[
\dim(Ker(u)) = 1
\]

---

\underline{Injectivité}

Non injective (noyau non nul)

---

\underline{Surjectivité}

Image = $\mathbb{R}^2$ ⇒ surjective

---

\underline{Conclusion}

\[
u \text{ est surjective mais non injective}
\]

---

%====================================================
\section{Corrigé — Méthodes}

%====================================================
\textbf{Exercice 6 — Construction}

\[
(x,y) = x e_1 + y e_2
\]

\[
u(x,y) = x u(e_1) + y u(e_2)
\]

\[
= x(1,2) + y(0,1)
\]

\[
= (x,\; 2x+y)
\]

---

%====================================================
\textbf{Exercice 7 — Matrice}

\[
u(1,0) = (1,2,0)
\]
\[
u(0,1) = (1,-1,1)
\]

---

\[
A =
\begin{pmatrix}
1 & 1 \\
2 & -1 \\
0 & 1
\end{pmatrix}
\]

---

%====================================================
\textbf{Exercice 8 — Noyau}

\[
\begin{pmatrix}
1 & 2 \\
2 & 4
\end{pmatrix}
\begin{pmatrix}x\\y\end{pmatrix} = 0
\]

\[
x+2y=0
\]

\[
x=-2y
\]

---

\[
Ker = Vect((-2,1))
\]

---

%====================================================
\section{Corrigé — Démonstrations}

%====================================================
\textbf{Exercice 10 — $u(0)=0$}

\[
u(0) = u(0+0) = u(0)+u(0)
\]

Donc :

\[
u(0)=0
\]

---

%====================================================
\textbf{Exercice 11 — Noyau SEV}

\underline{1. $0 \in Ker$}

\[
u(0)=0
\]

---

\underline{2. Stabilité addition}

Si $u(x)=0$ et $u(y)=0$ :

\[
u(x+y)=u(x)+u(y)=0
\]

---

\underline{3. Stabilité scalaire}

\[
u(\lambda x)=\lambda u(x)=0
\]

---

\underline{Conclusion}

$Ker(u)$ est un sous-espace vectoriel.
\section{Corrigé — Problèmes guidés (version corrigée)}

%====================================================
\textbf{Exercice 14 — Étude complète}

\[
u(x,y,z) = (x+y,\; y+z,\; x+z)
\]

---

\underline{1. Linéarité}

Chaque coordonnée est une combinaison linéaire de $x,y,z$, donc $u$ est linéaire.

---

\underline{2. Noyau}

On résout :
\[
\begin{cases}
x+y=0 \\
y+z=0 \\
x+z=0
\end{cases}
\]

De $x+y=0$ :
\[
y=-x
\]

De $y+z=0$ :
\[
z=x
\]

Dans la troisième équation :
\[
x+z = x+x = 2x = 0 \Rightarrow x=0
\]

Donc :
\[
y=0,\quad z=0
\]

\[
\boxed{Ker(u)=\{0\}}
\]

---

\underline{3. Conséquence}

Comme $u : \mathbb{R}^3 \to \mathbb{R}^3$ et que le noyau est nul :

\[
\boxed{u \text{ est injective}}
\]

---

\underline{4. Rang}

En dimension finie :

\[
\boxed{u \text{ injective } \Rightarrow u \text{ bijective}}
\]

Donc :

\[
\boxed{rg(u)=3}
\]

---

\underline{5. Vérification matricielle}

\[
A =
\begin{pmatrix}
1 & 1 & 0 \\
0 & 1 & 1 \\
1 & 0 & 1
\end{pmatrix}
\]

\[
\det(A) = 2 \neq 0
\]

---

\underline{Conclusion}

\[
\boxed{u \text{ est bijective}}
\]

\[
Im(u)=\mathbb{R}^3
\]

---

\underline{Interprétation géométrique}

L’application est un automorphisme de $\mathbb{R}^3$.

---

%====================================================
\textbf{Exercice 15 — Construction}

---

\underline{1. Unicité}

Une application linéaire est entièrement déterminée par l’image d’une base.

Donc l’application est unique.

---

\underline{2. Expression}

\[
(x,y,z)=x e_1 + y e_2 + z e_3
\]

\[
u(x,y,z)=x(1,0)+y(0,1)+z(1,1)
\]

\[
\boxed{u(x,y,z)=(x+z,\; y+z)}
\]

---

\underline{3. Matrice}

\[
A =
\begin{pmatrix}
1 & 0 & 1 \\
0 & 1 & 1
\end{pmatrix}
\]

---

\underline{4. Noyau}

\[
x+z=0,\quad y+z=0
\]

\[
z=-x,\quad y=x
\]

\[
\boxed{Ker(u)=Vect((1,1,-1))}
\]

---

\underline{5. Image}

Vecteurs colonnes :

\[
(1,0), (0,1), (1,1)
\]

Donc :

\[
\boxed{Im(u)=\mathbb{R}^2}
\]

---

\underline{6. Rang}

\[
\boxed{rg(u)=2}
\]

---

%====================================================
\textbf{Exercice 16 — Étude paramétrée}

\[
u_\lambda(x,y)=(x+\lambda y,\; y+\lambda x)
\]

---

\underline{1. Linéarité}

Combinaisons linéaires ⇒ linéaire.

---

\underline{2. Matrice}

\[
A =
\begin{pmatrix}
1 & \lambda \\
\lambda & 1
\end{pmatrix}
\]

---

\underline{3. Déterminant}

\[
\det(A)=1-\lambda^2
\]

---

\underline{4. Injectivité et surjectivité}

\[
\boxed{\lambda \neq \pm1 \Rightarrow u_\lambda \text{ bijective}}
\]

---

\underline{5. Cas $\lambda=1$}

\[
A =
\begin{pmatrix}
1 & 1 \\
1 & 1
\end{pmatrix}
\]

\[
x+y=0 \Rightarrow y=-x
\]

\[
\boxed{Ker = Vect((1,-1))}
\]

---

\underline{Rang}

\[
rg=1
\]

---

\underline{6. Cas $\lambda=-1$}

\[
A =
\begin{pmatrix}
1 & -1 \\
-1 & 1
\end{pmatrix}
\]

\[
x-y=0 \Rightarrow x=y
\]

\[
\boxed{Ker = Vect((1,1))}
\]

---

\underline{Rang}

\[
rg=1
\]

---

%====================================================
\textbf{Exercice 17 — Étude complète}

\[
u(x,y,z)=(x+y+z,\; x+2y+3z)
\]

---

\underline{1. Linéarité}

Évidente.

---

\underline{2. Images}

\[
u(1,0,0)=(1,1)
\]
\[
u(0,1,0)=(1,2)
\]
\[
u(0,0,1)=(1,3)
\]

---

\underline{3. Relation}

\[
(1,3)=2(1,2)-(1,1)
\]

Donc dépendance.

---

\underline{4. Base image}

\[
\boxed{Im(u)=Vect((1,1),(1,2))}
\]

---

\underline{5. Rang}

\[
\boxed{rg(u)=2}
\]

---

\underline{6. Noyau}

\[
\begin{cases}
x+y+z=0 \\
x+2y+3z=0
\end{cases}
\]

Soustraction :

\[
y+2z=0 \Rightarrow y=-2z
\]

\[
x -2z + z = 0 \Rightarrow x=z
\]

---

\[
\boxed{Ker = Vect((1,-2,1))}
\]

---

\underline{7. Vérification}

\[
3 = 1 + 2
\]

✔ OK

---

%====================================================
\section{Corrigé — ORAUX}

%====================================================
\textbf{Exercice 20 — Projecteur}

\[
u^2=u
\]

---

\[
x=(x-u(x)) + u(x)
\]

---

\underline{1.}

\[
u(x-u(x))=u(x)-u^2(x)=0
\]

Donc :

\[
x-u(x)\in Ker(u)
\]

---

\underline{2.}

\[
u(x)\in Im(u)
\]

---

\underline{3.}

\[
E = Ker(u) + Im(u)
\]

---

\underline{4. Directe}

Si $x \in Ker \cap Im$ :

\[
x=u(y)
\quad \text{et} \quad u(x)=0
\]

\[
u(x)=u^2(y)=u(y)=x
\Rightarrow x=0
\]

---

\[
\boxed{E=Ker(u)\oplus Im(u)}
\]

---

\underline{Conclusion}

$u$ est un projecteur.


\section{Corrigé — Problèmes avancés}

%====================================================
\textbf{Exercice 25 — Endomorphismes de rang 1}

Soit $u : E \to E$ de rang 1.

---

\underline{1. Image}

Par définition :

\[
rg(u)=1 \Rightarrow \dim(Im(u))=1
\]

Donc :

\[
\boxed{Im(u) \text{ est une droite vectorielle}}
\]

---

\underline{2. Colinéarité}

Soit $e \in Im(u)$ non nul.

Alors :

\[
Im(u)=Vect(e)
\]

Donc pour tout $x$ :

\[
u(x) \in Im(u) \Rightarrow u(x)=\lambda e
\]

---

\underline{3. Existence de $\varphi$}

On définit :

\[
\varphi(x)=\lambda
\]

Donc :

\[
\boxed{u(x)=\varphi(x)e}
\]

---

\underline{4. Linéarité de $\varphi$}

Soient $x,y$ :

\[
u(x+y)=\varphi(x+y)e
\]

Mais :

\[
u(x+y)=u(x)+u(y)=\varphi(x)e+\varphi(y)e
\]

Donc :

\[
\varphi(x+y)=\varphi(x)+\varphi(y)
\]

Même chose pour les scalaires ⇒ $\varphi$ linéaire.

---

\underline{5. Réciproque}

Si $u(x)=\varphi(x)e$ :

\[
Im(u) \subset Vect(e)
\]

Et si $\varphi \neq 0$ :

\[
Im(u)=Vect(e)
\]

Donc :

\[
rg(u)=1
\]

---

\underline{6. Noyau}

\[
u(x)=0 \Rightarrow \varphi(x)e=0 \Rightarrow \varphi(x)=0
\]

\[
\boxed{Ker(u)=Ker(\varphi)}
\]

---

%====================================================
\textbf{Exercice 26 — Projecteur}

\[
u^2=u
\]

---

\underline{1. Décomposition}

\[
x = (x-u(x)) + u(x)
\]

---

\underline{2.}

\[
u(x-u(x))=u(x)-u^2(x)=0
\]

Donc :

\[
x-u(x)\in Ker(u)
\]

---

\underline{3.}

\[
u(x)\in Im(u)
\]

---

\underline{4.}

\[
E = Ker(u) + Im(u)
\]

---

\underline{5. Directe}

Si $x \in Ker \cap Im$ :

\[
x=u(y), \quad u(x)=0
\]

\[
u(x)=u^2(y)=u(y)=x
\Rightarrow x=0
\]

---

\[
\boxed{E=Ker(u)\oplus Im(u)}
\]

---

\underline{Conclusion}

$u$ est le projecteur sur $Im(u)$ parallèlement à $Ker(u)$.

---

%====================================================
\textbf{Exercice 27 — Décomposition}

On suppose :

\[
E = Ker(u) \oplus F
\]

---

\underline{1. Injectivité sur $F$}

Si $x \in F$ et $u(x)=0$ :

\[
x \in Ker(u) \cap F
\]

Donc :

\[
x=0
\]

---

\underline{2. Image}

\[
Im(u)=u(E)=u(F)
\]

---

\underline{3. Dimension}

\[
\dim(Im(u))=\dim(F)
\]

---

\underline{4. Théorème du rang}

\[
\dim(E)=\dim(Ker(u))+\dim(F)
\]

Donc :

\[
\boxed{\dim(E)=\dim(Ker(u))+rg(u)}
\]

---

%====================================================
\textbf{Exercice 28 — Involution}

\[
u^2=Id
\]

---

\underline{1.}

\[
(u-Id)(u+Id)=u^2-Id=0
\]

---

\underline{2.}

\[
Im(u-Id) \subset Ker(u+Id)
\]

---

\underline{3. Décomposition}

Soit $x$ :

\[
x = \frac{1}{2}(x+u(x)) + \frac{1}{2}(x-u(x))
\]

---

On montre :

\[
\frac{1}{2}(x+u(x)) \in Ker(u-Id)
\]
\[
\frac{1}{2}(x-u(x)) \in Ker(u+Id)
\]

---

\underline{4. Somme directe}

Intersection triviale ⇒ somme directe.

---

\[
\boxed{E = Ker(u-Id) \oplus Ker(u+Id)}
\]

---

%====================================================
\section{Corrigé — ORAUX}

%====================================================
\textbf{Exercice 29 — Théorème du rang}

---

Soit $(e_1,\dots,e_k)$ base du noyau.

On complète en base $(e_1,\dots,e_n)$.

---

\underline{1.}

Les vecteurs $u(e_1),\dots,u(e_k)=0$.

---

\underline{2.}

On montre que $u(e_{k+1}),\dots,u(e_n)$ sont libres :

Supposons :

\[
\sum \lambda_i u(e_i)=0
\Rightarrow u\left(\sum \lambda_i e_i\right)=0
\]

Donc combinaison dans le noyau ⇒ coefficients nuls.

---

\underline{3.}

Ils engendrent l’image.

---

\underline{Conclusion}

\[
\boxed{\dim(E)=\dim(Ker(u))+\dim(Im(u))}
\]

---

%====================================================
\textbf{Exercice 30 — Nilpotent}

\[
u^2=0
\]

---

\underline{1.}

\[
u(u(x))=0 \Rightarrow u(x)\in Ker(u)
\]

Donc :

\[
Im(u)\subset Ker(u)
\]

---

\underline{2.}

\[
rg(u)\leq \dim(Ker(u))
\]

---

\underline{3.}

\[
\dim(E)=\dim(Ker(u))+rg(u)
\]

Donc :

\[
rg(u)\leq \frac{\dim(E)}{2}
\]

---

%====================================================
\textbf{Exercice 31 — Isomorphisme}

---

Si injective :

\[
\dim(Im(u))=\dim(E)=\dim(F)
\]

Donc :

\[
Im(u)=F
\]

---

Donc bijective.

---

Réciproque identique.

---

\[
\boxed{u \text{ bijective } \Longleftrightarrow \text{ injective ou surjective}}
\]

---

%====================================================
\textbf{Exercice 32 — Structure}

---

\underline{1.}

On complète une base du noyau ⇒ existence de $G$.

---

\underline{2.}

Restriction injective comme précédemment.

---

\underline{3.}

\[
Im(u)=u(G)
\]

---

\underline{4.}

Application bijective entre $G$ et $Im(u)$.

---

\underline{Conclusion}

\[
\boxed{u \text{ induit un isomorphisme entre } G \text{ et } Im(u)}
\]
\end{document}